% ---------------------------------------------------------------------
%  Chapter 8 : The Mean-Field Saddle and the DAE Solver
% ---------------------------------------------------------------------
\chapter{The Mean-Field Saddle and the DAE Solver}\label{ch:mean-field}

Every loop expansion needs a point to expand \emph{about}. In the \msrjd{}
formalism that point is the \term{mean-field saddle}: the deterministic steady
state of the system with all fluctuations switched off. It is the firing rate
$n^{*}$, the membrane voltage $v^{*}$, the variance-suppressed equilibrium value
$x^{*}$ --- whatever the deterministic field settles to when the noise is muted.
The whole perturbative series (the diagrams) is a systematic correction
\emph{around} that background, so before a single diagram can be evaluated the
engine must find the saddle and turn it into numbers.

This chapter is about the subsystem that does exactly that. Its single job is
narrow and entirely concrete: \textbf{given a fully specified model and a
dictionary of numeric parameter values, find the saddle and package it into a
substitution dictionary that the rest of the pipeline plugs into its symbolic
expressions}. That dictionary is called \code{num\_params}, and it is the only
thing this subsystem hands downstream. The chapter walks the two solvers that
produce it --- one for theories that declare their steady state as a
self-consistency closure, one for theories that declare it as a
differential-algebraic system --- and then the linear-stability machinery that
decides whether a given saddle is even a legitimate place to expand.

The code lives in two files. The legacy \emph{iteration} solver is
\file{pipeline/\_mean\_field.py} (553 lines); the \emph{DAE multi-root} solver
plus linear stability is \file{pipeline/\_mean\_field\_dae.py} (977 lines). They
are, by design, completely independent of each other, and which one runs is
decided by a single line in \file{pipeline/compute.py}.

This is phase \code{[3/7]} of the seven-phase trace you met in
Chapter~\ref{ch:quickstart} (\code{[3/7] Solve MF self-consistency...}). It sits
after the propagator has been built symbolically and before the propagator's
poles and residues are computed numerically --- because those poles need the
saddle's numbers to exist.

\section{Why a saddle at all}

\subsection{The background expansion in MSR-JD}

Recall the \msrjd{} picture. A stochastic dynamical system --- a Langevin
equation, a master equation, a spiking network --- is rewritten as a generating
functional
\begin{equation}
  Z[J] \;=\; \int \mathcal{D}[\phi]\,\mathcal{D}[\tilde\phi]\;
    \exp\!\big(-S[\phi,\tilde\phi] + \text{sources}\big),
  \label{eq:mf-genfunc}
\end{equation}
where $\phi$ is the physical field (a rate, a voltage, a density) and
$\tilde\phi$ is the conjugate \term{response field}. The action $S$ is the object
the entire pipeline manipulates.

Perturbation theory expands the physical field about a background,
$\phi = \phi^{*} + \delta\phi$, and the background $\phi^{*}$ is chosen so that
the action has \emph{no linear term} in the fluctuation $\delta\phi$. That
``vanishing linear term'' condition is not an arbitrary convenience: it is
exactly the classical equation of motion, i.e.\ the deterministic steady state.
A linear term in $\delta\phi$ would be a tadpole --- a one-point function that
should vanish on a correctly-chosen background --- and leaving it in would make
the diagrammatic bookkeeping wrong. So the saddle is the \emph{tree-level}
solution, and the loop corrections (the diagrams) are the systematic corrections
to it.

\begin{defn}
The \textbf{mean-field saddle} $\phi^{*}$ is the field configuration that makes
the action's first variation vanish, $\delta S/\delta\phi|_{\phi^{*}} = 0$. It is
the deterministic steady state of the dynamics. The loop expansion is an
expansion in fluctuations $\delta\phi = \phi - \phi^{*}$ \emph{around} this point.
\end{defn}

\subsection{The canonical self-consistency problem}

For a rate network the saddle is a self-consistency problem. The textbook neural
example, which the legacy solver is built around, is
\begin{align}
  v^{*}_i &= E_i + \sum_j w_{ij}\, g_{ij}\!\star n^{*}_j,
    && \text{(voltage $=$ drive $+$ recurrent input)} \label{eq:mf-volt}\\
  n^{*}_i &= \phi_i(v^{*}_i),
    && \text{(rate $=$ transfer function of voltage).} \label{eq:mf-rate}
\end{align}
Here $g_{ij}$ is a synaptic kernel, $E_i$ an external drive, $w_{ij}$ a coupling
weight, and $\phi_i$ a (possibly nonlinear) \term{transfer function} from input
to output rate. The symbol $\star$ is convolution in time.

One feature of \eqref{eq:mf-volt} is load-bearing throughout the code: at the
saddle every kernel \emph{integrates to one}. Kernels are normalized,
$\int g(t)\,\dd t = 1$, so the convolution $g\star n^{*}$ of a \emph{constant}
$n^{*}_j$ is just $n^{*}_j$ --- a stationary mean field sees only a kernel's total
weight, never its shape. This is precisely why both solvers substitute every
kernel symbol $\to 1$ before evaluating anything (you will see this as
\code{g\_to\_one} and \code{kernel\_to\_one} below).

Substituting the voltage equation \eqref{eq:mf-volt} into the rate equation
\eqref{eq:mf-rate} closes the problem in the rates alone:
\begin{equation}
  n^{*}_i \;=\; \phi_i\!\Big(E_i + \sum_j w_{ij}\, n^{*}_j\Big)
    \;\equiv\; T_i(n^{*}).
  \label{eq:mf-closure}
\end{equation}
The legacy solver finds $n^{*}$ by driving the residual $n^{*}_i - T_i(n^{*})$ to
zero. This is the split that pervades \file{\_mean\_field.py}: the rates $n^{*}$
are the \term{iteration saddle} (the variables the root finder actually
iterates), and the voltages $v^{*}$ are the \term{compound saddle} (computed from
the rates on the fly each step).

\subsection{Where it sits in the pipeline}

The subsystem is fed by \code{TheoryBuilder.build()} (the \code{model} dict plus a
compiled \code{FieldTheory}) and the user's numeric parameters, and it hands its
\code{num\_params} dict to the propagator's pole finder. Schematically:

\begin{lstlisting}[style=console]
   TheoryBuilder.build()              (pipeline/theory.py)
        |  model dict + compiled FieldTheory ft
        v
   compute_cumulants(...)             (pipeline/compute.py)
        |  [1] expand action -> vertices / sources
        |  [2] build_propagator -> K_ft, G_ft   (symbolic, in omega)
        |
        +--[3] MEAN-FIELD SOLVE  <----------  THIS SUBSYSTEM
        |        if model['equations']: solve_mean_field_dae_compat(...)
        |        else:                   solve_mean_field(...)
        |        -> mf['num_params']   {SR symbol -> float}
        |
        v
   compute_poles_and_residues(prop, num_params)   (pipeline/_propagator.py)
        |  substitute num_params into K_ft, find retarded poles, residues
        v
   Phase J / time-domain machinery -> connected cumulant C(tau)
\end{lstlisting}

What feeds it:
\begin{itemize}
  \item \code{ft} --- a compiled \code{FieldTheory}. The solver touches only
        \code{ft.\_ns} (the Sage symbol namespace) and \code{ft.taylor\_order}
        (how many transfer-function derivatives to precompute). In the DAE path
        \code{ft} is used \emph{only} by the compatibility wrapper, to map saddle
        names back onto the \code{ft.\_ns} symbol arrays.
  \item \code{model} --- the declaration dict from \code{build()}. It carries
        \code{parameters}, \code{kernels}, \code{populations},
        \code{physical\_fields}, \code{functions}, \code{phi\_concrete},
        \code{mf\_bg\_conditions}, \code{iteration\_saddles}, \code{equations},
        and \code{stability\_analysis}.
  \item \code{fundamental} --- the user's concrete numeric parameter values, keyed
        by parameter \emph{name} (e.g.\ \code{\{'E': [...], 'w': [[...]], 'tau': 1.0\}}).
\end{itemize}

What consumes its output: the \emph{primary} consumer is
\code{compute\_poles\_and\_residues(prop, num\_params)}
(\file{pipeline/\_propagator.py:851}), which substitutes \code{num\_params} into
the symbolic kernel $K_{\mathrm{ft}}(\omega)$ to get a numeric matrix-valued
rational function of $\omega$, then finds its retarded poles and residue
matrices. A second consumer, \code{\_solve\_mf\_at\_saddle}
(\file{pipeline/\_precompute.py:200}), re-runs the solver to verify the action's
linear term vanishes at the saddle. The handoff from the saddle to the
perturbative expansion is, in both cases, \emph{purely substitutional}.

\section{Two solvers, one routing decision}

A model declares its steady state in one of two styles, and the framework picks a
solver to match. The decision is a single truthiness test on
\code{model['equations']}:

\begin{lstlisting}
# pipeline/compute.py:364
if model.get('equations'):
    mf = solve_mean_field_dae_compat(
        ft, model, fundamental,
        fixed_point_index=fixed_point_index,
        n_starts=mf_dae_n_starts,
        seed_box=mf_dae_seed_box,
        verbose=verbose,
    )
else:
    mf = solve_mean_field(ft, model, fundamental, verbose=verbose)
num_params = mf['num_params']
\end{lstlisting}

The two declaration styles, and the solvers they route to:

\begin{description}
  \item[Self-consistency closure $\to$ legacy iteration solver.] A theory whose
        saddle is described as a closure ($n^{*}_i = \phi(v^{*}_i)$ with $v^{*}_i$
        an algebraic function of the $n^{*}_j$) goes to
        \code{solve\_mean\_field} in \file{\_mean\_field.py}. It iterates only the
        rate-like saddle variables and treats the rest (voltages) as compound
        expressions evaluated on the fly. It has a single-population branch and a
        heterogeneous (\code{E}/\code{I}) branch.
  \item[Differential-algebraic system $\to$ DAE multi-root solver.] A theory that
        declares its dynamics directly, one equation per field, via
        \code{TheoryBuilder.equation(lhs, rhs, population)}
        (\file{pipeline/theory.py:1087}), goes to
        \code{solve\_mean\_field\_dae\_compat} in \file{\_mean\_field\_dae.py}. It
        sets the time-derivative operator $\code{Dt}\to 0$ on every equation,
        solves the resulting algebraic root problem with many random restarts,
        deduplicates the roots, optionally classifies each root's linear
        stability, and selects one root by index.
\end{description}

\begin{defn}
\textbf{DAE} stands for \emph{differential-algebraic equation}: a system that
mixes genuine differential equations (those carrying the time-derivative operator
\code{Dt}, such as a voltage relaxation $(\tau\,\code{Dt}+1)v=\dots$) with purely
algebraic constraints (those without \code{Dt}, such as a rate readout
$n=\phi(v)$). At the saddle the differential equations collapse to algebraic ones
($\code{Dt}\to 0$), but the differential structure is \emph{remembered} and
reused for stability analysis.
\end{defn}

\begin{gotcha}
The routing is \emph{purely} the truthiness of \code{model.get('equations')}
(\file{compute.py:364}, and identically at \file{\_precompute.py:204}). A theory
that declares both \code{.equation(...)} \emph{and} \code{mf\_bg\_conditions} will
silently take the DAE path. The two modules are, in the words of the source,
``INTENTIONALLY decoupled'' (\file{\_mean\_field\_dae.py:16}) --- there is no
shared code and no negotiation between them.
\end{gotcha}

\section{The legacy iteration solver}

\code{solve\_mean\_field(ft, model, fundamental, verbose=True)}
(\file{\_mean\_field.py:14}) is the entry point. It immediately routes to the
heterogeneous branch if the model declares populations:

\begin{lstlisting}
# pipeline/_mean_field.py:40
if model.get('populations'):
    return _solve_mean_field_hetero(ft, model, fundamental, verbose)
\end{lstlisting}

We will follow the single-population path first, then note what the
heterogeneous branch adds. The solver returns a five-key dict:
\code{nstar\_vals} (list of floats), \code{vstar\_vals} (list),
\code{phi\_deriv\_vals} (a \code{\{(dk, i): float\}} map of saddle derivatives),
\code{num\_params} (the deliverable), and \code{param\_subs} (the raw user
parameters expanded to indexed symbols).

\subsection{Step 1 --- expand the parameters into Sage symbols}

The first job is to turn the user's parameter values, keyed by name, into a dict
keyed by the \emph{indexed Sage symbols} that actually appear inside the
propagator. A matrix parameter \code{w} becomes \code{w11}, \code{w12}, \dots; a
vector becomes \code{E1}, \code{E2}, \dots; a scalar stays bare. The axis sizes
come from the population-size map \code{ns.\_pop\_size}, but the \emph{shape of
the user's value} is the source of truth for whether something is scalar, vector,
or matrix:

\begin{lstlisting}
# pipeline/_mean_field.py:62-87  (abridged)
ib = pspec.get('indexed_by')
if ib:
    if len(ib) == 2:
        n_rows = pop_size_map.get(ib[0], len(ns.pop))
        n_cols = pop_size_map.get(ib[1], len(ns.pop))
        for i in range(n_rows):
            for j in range(n_cols):
                param_subs[SR.var(f'{pname}{i+1}{j+1}')] = val[i][j]
    elif len(ib) == 1:
        ...
        param_subs[SR.var(f'{pname}{i+1}')] = val[i]
    ...
else:
    param_subs[SR.var(pname)] = val
\end{lstlisting}

\code{SR.var('w12')} returns (and globally registers) the symbolic variable named
\code{w12}; calling it again anywhere with the same name returns the \emph{same}
symbol object. That identity is exactly why a substitution dict keyed by
\code{SR.var('w12')} matches the \code{w12} that the propagator built
independently. (Section~\ref{sec:mf-tools} explains Sage's symbolic ring from the
ground up.)

\subsection{Step 2 --- the transfer-function derivatives, symbolically}

The diagrammatic vertices are obtained by Taylor-expanding the action about the
saddle. A transfer function $\phi$ that is nonlinear in the field contributes
vertices whose couplings are the \emph{derivatives of $\phi$ evaluated at the
saddle}:
\begin{equation}
  \phi(v^{*} + \delta v) = \phi(v^{*}) + \phi'(v^{*})\,\delta v
    + \tfrac{1}{2}\phi''(v^{*})\,\delta v^{2}
    + \tfrac{1}{6}\phi'''(v^{*})\,\delta v^{3} + \cdots
  \label{eq:mf-taylor}
\end{equation}
The pipeline names these $\code{phi0\_i} = \phi(v^{*}_i)$ (which equals the rate
$n^{*}_i$), $\code{phi1\_i} = \phi'(v^{*}_i)$, $\code{phi2\_i} = \phi''(v^{*}_i)$,
and so on up to \code{ft.taylor\_order}. The solver builds them by symbolic
differentiation of the \emph{concrete} $\phi$ expression and then evaluating at
the numeric $v^{*}_i$:

\begin{lstlisting}
# pipeline/_mean_field.py:90-102
v_sym = SR.var('_v_mf_')
phi_derivs = {}
for i in ns.pop:
    phi_expr = model['phi_concrete'](ns, i, v_sym)
    phi_derivs[i] = {}
    for dk in range(taylor_order + 1):
        if dk == 0:
            phi_derivs[i][dk] = phi_expr
        else:
            phi_derivs[i][dk] = diff(phi_expr, v_sym, dk)

def phi_num(i, v_val):
    return float(phi_derivs[i][0].subs(param_subs).subs({v_sym: v_val}))
\end{lstlisting}

Here \code{model['phi\_concrete'](ns, i, v\_sym)} returns the concrete symbolic
expression for $\phi$ on a dummy variable \code{\_v\_mf\_}, and
\code{diff(phi\_expr, v\_sym, dk)} takes its $dk$-th derivative. \code{phi\_num}
substitutes the parameters and the numeric voltage, then coerces to a Python
\code{float}.

\begin{defn}
The symbol $\code{phik\_i}$ denotes the $k$-th derivative of the transfer
function at the saddle of population $i$: $\code{phik\_i} = \phi^{(k)}(v^{*}_i)$.
The zeroth one is special, $\code{phi0\_i} = \phi(v^{*}_i) = n^{*}_i$. These are
the numeric vertex couplings the diagrams need.
\end{defn}

\subsection{Step 3 --- the symbolic voltage and the iteration target}

The voltage saddle is read \emph{symbolically} from the model rather than
hardcoded, so that any model whose voltage carries extra terms (feedforward
drives, additional couplings) is picked up automatically. The kernel symbol is
forced to $1$ and the parameters baked in:

\begin{lstlisting}
# pipeline/_mean_field.py:105-118
vstar_subs_dict = model.get('mf_bg_conditions', lambda ns: {})(ns)
vstar_sym = {i: SR(vstar_subs_dict.get(ns.vstar[i], ns.E[i])) for i in ns.pop}
g_to_one = {ns.g: SR(1)}
vstar_baked = {i: SR(vstar_sym[i]).subs(g_to_one).subs(param_subs) for i in ns.pop}

def vstar_num(i, nstar_vec):
    sub = {ns.nstar[j]: float(nstar_vec[j]) for j in ns.pop}
    return float(vstar_baked[i].subs(sub))
\end{lstlisting}

\code{mf\_bg\_conditions} is a callable on the namespace that returns a dict
mapping each saddle LHS symbol to its RHS expression. The line
\code{vstar\_subs\_dict.get(ns.vstar[i], ns.E[i])} reads the declared RHS for
$v^{*}_i$, defaulting to the bare drive \code{ns.E[i]} if the model declares
none. \code{g\_to\_one} is the kernel-normalization substitution discussed above
(\code{ns.g} $\to 1$).

The iteration target for $n^{*}$ is likewise taken from the model when declared,
falling back to the legacy hardcode $n^{*} = \phi(v^{*})$:

\begin{lstlisting}
# pipeline/_mean_field.py:124-141  (abridged)
nstar_rhs_baked = None
if hasattr(ns, 'nstar') and ns.nstar[0] in vstar_subs_dict:
    nstar_rhs_baked = {i: SR(vstar_subs_dict[ns.nstar[i]])
                            .subs(g_to_one).subs(param_subs) for i in ns.pop}

def nstar_target(i, nstar_vec):
    if nstar_rhs_baked is None:
        return phi_num(i, vstar_num(i, nstar_vec))   # legacy: n* = phi(v*)
    ...
\end{lstlisting}

\subsection{Step 4 --- multi-start \code{fsolve} with a singularity guard}

The self-consistency residual is $n^{*}_i - T_i(n^{*})$, and it is wrapped so that
symbolic singularities (a spike-reset RHS like $1/(1-n^{*})$ evaluated near
$n^{*}=1$) become a large \emph{finite} penalty rather than raising --- which lets
the root finder back away and keep searching:

\begin{lstlisting}
# pipeline/_mean_field.py:149-179  (abridged)
def mf_residual(nstar_vec):
    try:
        return [nstar_vec[i] - nstar_target(i, nstar_vec) for i in ns.pop]
    except (ValueError, ZeroDivisionError):
        return [1e10] * len(ns.pop)

initial_guesses = [[0.1]*npop, [0.5]*npop, [1.0]*npop, [0.01]*npop]
sol = None
for x0 in initial_guesses:
    try:
        attempt = fsolve(mf_residual, x0, full_output=True)
    except (ValueError, ZeroDivisionError):
        continue
    sol = attempt
    if attempt[2] == 1:        # ier == 1  =>  converged
        break
\end{lstlisting}

\code{fsolve} (SciPy's wrapper around MINPACK's Powell hybrid method, see
Section~\ref{sec:mf-tools}) returns a 4-tuple under \code{full\_output=True};
element \code{[2]} is the integer convergence flag \code{ier}, and \code{ier == 1}
means clean convergence. The solver tries the guesses $0.1, 0.5, 1.0, 0.01$ in
order and keeps the first that converges. The small leading guesses are
deliberate: they are safe for theories whose saddle sits below a pole at larger
$n^{*}$.

After the solve, explicit finiteness checks guard against a root finder that dove
into a near-pole and returned garbage:

\begin{lstlisting}
# pipeline/_mean_field.py:193-197
if not all(__import__('math').isfinite(x) for x in nstar_vals):
    raise RuntimeError(
        f'solve_mean_field: fsolve returned non-finite n* = '
        f'{nstar_vals!r}.  The mf_eq probably has a singularity '
        f'within the iteration basin; check parameter values.')
\end{lstlisting}

\begin{gotcha}
The legacy solver has \emph{no} multi-root machinery. It tries the four guesses
in order and returns whatever the first converging one lands on. A theory with a
saddle far outside those basins may converge to the \emph{wrong} fixed point or
fail outright. If you need to discover multiple fixed points or pick a specific
branch, that is precisely what the DAE solver is for.
\end{gotcha}

\subsection{Step 5 --- assemble \code{num\_params} and the \code{phi0} rescue}

With the saddle in hand, the voltages and $\phi$-derivatives are evaluated at the
fixed point and everything is poured into \code{num\_params}: the raw parameters,
the saddle values \code{ns.nstar[i]}, \code{ns.vstar[i]}, the GTaS rate
\code{ns.mstar[i]} (guarded behind a presence check), and the numeric derivative
symbols \code{phi1\_i}, \code{phi2\_i}, \dots

\begin{lstlisting}
# pipeline/_mean_field.py:234-258  (abridged)
num_params = dict(param_subs)
for i in ns.pop:
    num_params[ns.nstar[i]] = float(nstar_vals[i])
    num_params[ns.vstar[i]] = float(vstar_vals[i])
...
for i in ns.pop:
    for dk in range(1, taylor_order + 1):
        sym = SR.var(f'phi{dk}_{i+1}')
        if (dk, i) in phi_deriv_vals:
            num_params[sym] = phi_deriv_vals[(dk, i)]
\end{lstlisting}

The last act is the most easily overlooked, and the source comments it heavily.
Template-built theories keep $\phi$ as a \emph{formal} function symbol, so the
$(2,0)$ Poisson-noise sector of the action carries the formal symbol
\code{phi0\_i} (the mean rate $n^{*}_i$) rather than a number. The per-derivative
loop above only resolves \code{phi1\_i} upward. If \code{phi0\_i} is left
symbolic, the noise-source coefficient $-\tfrac12\,\code{phi0\_i}$ never becomes
numeric --- and \emph{the entire correlator collapses to zero}. The fix is to
evaluate the remaining \code{mf\_bg\_conditions} entries (notably the identity
$\code{phi0\_i}\to n^{*}_i$) at the solved saddle:

\begin{lstlisting}
# pipeline/_mean_field.py:272-279
for _lhs, _rhs in vstar_subs_dict.items():
    if _lhs in num_params:
        continue          # saddle vars already resolved above
    try:
        num_params[_lhs] = float(SR(_rhs).subs(g_to_one).subs(num_params))
    except (TypeError, ValueError):
        pass              # still symbolic -- leave it out
\end{lstlisting}

\begin{gotcha}
If a new template forgets to map \code{phi0\_i} $\to$ \code{nstar\_i} in its
\code{mf\_bg\_conditions}, the symbol stays formal, the noise vertex
$-\tfrac12\,\code{phi0\_i}$ never goes numeric, and the correlator silently
evaluates to zero. This is a quiet failure mode, not a crash --- watch for an
all-zero $C(\tau)$ when bringing up a new template-built theory.
\end{gotcha}

\subsection{The heterogeneous branch}

\code{\_solve\_mean\_field\_hetero} (\file{\_mean\_field.py:291}) handles multiple
populations (e.g.\ excitatory and inhibitory), each with its own $\phi$ and its
own saddle arrays. It generalizes the single-pop path along three axes. First, it
identifies which parameters are saddles (those with \code{mean\_field=True}) and
records the population each ranges over (\file{\_mean\_field.py:338-351}). Second,
it classifies saddles as iteration vs.\ compound, preferring the authoritative
\code{model['iteration\_saddles']} list computed at build time and falling back to
``all are iteration'' only if it is missing (\file{\_mean\_field.py:367-375}).
Third, it builds a \emph{flat} fsolve vector that concatenates all
iteration-saddle elements, with helpers \code{\_unflatten} and
\code{\_eval\_compound} to convert between the flat vector and per-saddle arrays
and to evaluate the compound (voltage) saddles from the iteration (rate) values
each step. The kernel-to-one normalization here maps \emph{every} kernel symbol
on the namespace --- scalar, vector, and matrix --- to $1$
(\file{\_mean\_field.py:387-401}).

\begin{note}
The heterogeneous path returns \code{phi\_deriv\_vals: \{\}} (empty). In the
multi-population case the propagator builds the $\phi$ derivatives symbolically
rather than relying on this dict. The same is true of the entire DAE path. If a
consumer ever needs explicit saddle derivatives down those routes, it would have
to populate this dict --- the code does not do so today.
\end{note}

\section{The DAE formulation}

The DAE path declares the dynamics directly, one equation per field, through
\code{TheoryBuilder.equation(lhs, rhs, population)}. The neural example becomes
\begin{align}
  (\tau_i\,\code{Dt} + 1)\,v_i &= E_i + \textstyle\sum_j w_{ij}\,n_j
    && \text{(differential --- carries \code{Dt})}, \\
  n_i &= \phi_i(v_i)
    && \text{(algebraic --- no \code{Dt})}.
\end{align}
The symbol \code{Dt} is a formal stand-in for the time-derivative operator
$\dd/\dd t$. At the stationary mean-field saddle everything is time-independent,
so $\code{Dt}\to 0$ and the system becomes purely algebraic:
\begin{equation}
  v_i = E_i + \sum_j w_{ij}\,n_j,
  \qquad
  n_i = \phi_i(v_i).
\end{equation}
The solver hard-substitutes $\code{Dt}=0$ (and, for spatial theories,
$\code{Laplacian}=0$, since a spatially-uniform background has zero Laplacian)
into every equation and finds the joint root $(v^{*}, n^{*})$. The residual vector
is $F_k = \text{LHS}_k - \text{RHS}_k$, and the root satisfies $F(x^{*}) = 0$.

\subsection{Evaluating the residual via sandboxed \code{eval}}

This is the design choice that defines the DAE solver: it does \emph{not} parse
the user's equation text into a syntax tree. It \code{eval}s the text directly,
in a carefully constructed namespace. The residual builder
\code{\_build\_residual} (\file{\_mean\_field\_dae.py:236}) first expands each
declared equation over its population's index range into a list of single-scalar
residuals, then returns a closure \code{R(x)} that packs the flat state vector
into per-variable NumPy arrays and evaluates each residual:

\begin{lstlisting}
# pipeline/_mean_field_dae.py:268-307  (abridged)
def R(x: np.ndarray) -> np.ndarray:
    state = {var: np.asarray(x[start:end], dtype=float)
             for var, (start, end) in slices.items()}
    ns = {
        **params_np, **state, **phi_callables, **index_sets, **_MATH_NS,
        'Dt':           0,
        'Laplacian':    0,
        'sum':          sum,
        '__builtins__': {},
    }
    out = np.empty(len(expanded), dtype=float)
    for k, (lhs_text, rhs_text, i, pop) in enumerate(expanded):
        ns['i'] = i
        try:
            lhs = eval(lhs_text, ns)
            rhs = eval(rhs_text, ns)
            out[k] = float(lhs - rhs)
        except (ValueError, ZeroDivisionError, OverflowError, FloatingPointError):
            out[k] = 1e10
    return out
\end{lstlisting}

Several details here are load-bearing. The math functions in \code{\_MATH\_NS}
(\file{\_mean\_field\_dae.py:33-38}) are drawn from NumPy (\code{np.tanh},
\code{np.exp}, \dots) so the expression evaluates correctly over both scalars and
arrays. \code{Dt} and \code{Laplacian} are bound to $0$.
\code{'\_\_builtins\_\_': \{\}} sandboxes the evaluation so that only the
explicitly-bound names are visible. Exceptions from singular regions become the
$10^{10}$ penalty that steers the solver away.

\begin{gotcha}
The namespace is passed to \code{eval} as the \emph{globals} argument, not
locals. Python's comprehension-scoping rule (PEP 3104) runs a generator
expression such as \code{sum(w[i,j]*n[j] for j in E)} in its own function scope,
which sees the \code{eval}-globals but \emph{not} the \code{eval}-locals. Passing
the namespace as locals would make the genexpr raise \code{NameError} on
\code{w}, \code{n}, \code{E}. This subtlety is documented in the source at
\file{\_mean\_field\_dae.py:273-279}.
\end{gotcha}

\begin{gotcha}
Because \code{'\_\_builtins\_\_'} is empty, \emph{only} the names the solver
explicitly binds are available (parameters, state variables, $\phi$ callables,
index sets, the \code{\_MATH\_NS} functions, \code{Dt}, \code{Laplacian},
\code{sum}). A user equation that reaches for any other Python builtin --- say
\code{min}, \code{max}, or \code{len} --- will raise, and that exception is
swallowed into the $10^{10}$ penalty, which can masquerade as a failure to
converge. If a DAE theory mysteriously finds no root, an unsupported builtin in
the equation text is a prime suspect.
\end{gotcha}

\begin{note}
\code{\_sage\_to\_python} (\file{\_mean\_field\_dae.py:41}) translates the power
operator \code{\^{}} to \code{**} before any \code{eval}. Sage preparses
\code{\^{}} as exponentiation, but in plain Python \code{\^{}} is bitwise XOR, so
\code{eps*x[i]\^{}3} would raise \code{TypeError: ufunc 'bitwise\_xor' not
supported} on a NumPy float. The replacement is unconditional and safe, since XOR
is never meaningful in a mean-field equation over reals.
\end{note}

\subsection{The transfer functions as numeric callables}

The $\phi$ functions referenced in the equation RHSs are turned into numeric
callables by \code{\_make\_phi\_callables} (\file{\_mean\_field\_dae.py:147}).
Each declared single-argument function becomes a per-population closure that
\code{eval}s the function's expression text in the same sandboxed style. They are
wrapped in a small helper class, \code{\_PhiCallableList}
(\file{\_mean\_field\_dae.py:114}), that supports both indexed access
\code{phi[i](v)} and --- when there is exactly one population position --- a bare
scalar call \code{phi(v)}:

\begin{lstlisting}
# pipeline/_mean_field_dae.py:134-144
def __call__(self, *args, **kwargs):
    if len(self._callables) != 1:
        raise TypeError(
            f'bare call ``f(...)`` requires exactly one population '
            f'position; got {len(self._callables)}.  Use ``f[i](...)`` '
            f'with an explicit index instead.')
    return self._callables[0](*args, **kwargs)
\end{lstlisting}

This mirrors the action-side \code{\_IndexedFormalFunction} convention, so a user
can write \code{phi(v)} in a single-population theory and \code{phi[i](v)} in a
multi-population one.

\begin{gotcha}
\code{\_make\_phi\_callables} only handles \emph{single-argument} functions
(\file{\_mean\_field\_dae.py:174-175}: \code{if len(args\_text) != 1: continue}).
Multi-argument functions --- e.g.\ CGF-style cumulant generating functions --- are
silently skipped here. The DAE solver only needs single-argument transfer
functions for the saddle, so this is by design, but a multi-argument function
referenced in an equation RHS will surface as an unbound name (and hence the
swallowed penalty).
\end{gotcha}

\section{Finding all fixed points}

Nonlinear self-consistency equations generically have \emph{several} solutions:
bistability, up/down states, multiple firing-rate fixed points. The DAE solver
embraces this with a multi-start strategy. The entry point is
\code{solve\_mean\_field\_dae} (\file{\_mean\_field\_dae.py:390}).

\subsection{Seeds, the seed box, and domains}

It draws \code{n\_starts} random initial guesses (default $64$) from a
per-variable box. The box is domain-aware: a variable declared
\code{domain='positive'} is sampled from $[0,\,5s]$, anything else from
$[-3s,\,3s]$, where the scale $s$ is the largest absolute parameter value
(default $1.0$). The default domain, when a field declares none, is
\code{'positive'} --- the Hawkes-like assumption that rates are non-negative:

\begin{lstlisting}
# pipeline/_mean_field_dae.py:336-342
domain = (field_specs.get(var, {}).get('domain') or 'positive')
if domain == 'positive':
    boxes[var] = (0.0, 5.0 * scale)
else:
    boxes[var] = (-3.0 * scale, 3.0 * scale)
\end{lstlisting}

The RNG is the modern \code{np.random.default\_rng(rng\_seed)}
(\file{\_mean\_field\_dae.py:475}), seeded for reproducibility, so the same theory
and parameters always sample the same seeds. A caller can override the box for
any variable via the \code{seed\_box} argument.

\subsection{Multi-start Newton with a re-checked residual}

Each seed is handed to SciPy's \code{root} with the \code{hybr} method (the same
Powell hybrid as \code{fsolve}, but returning a richer result object). The solver
keeps only roots that survive three filters: SciPy's own success flag, a
\emph{re-evaluated} residual below $10^{-7}$, and finiteness:

\begin{lstlisting}
# pipeline/_mean_field_dae.py:481-499  (abridged)
for s in range(n_starts):
    x0 = seeds[s]
    try:
        sol = _scipy_root(R, x0, method='hybr')
    except (ValueError, ZeroDivisionError, OverflowError, FloatingPointError):
        continue
    if not sol.success:
        continue
    resid = R(sol.x)
    if np.max(np.abs(resid)) > 1e-7:
        continue
    if not np.all(np.isfinite(sol.x)):
        continue
    n_converged += 1
    raw_roots.append(np.asarray(sol.x, dtype=float))
\end{lstlisting}

\begin{gotcha}
SciPy's \code{sol.success} is not trusted. The solver re-evaluates \code{R(sol.x)}
and rejects any root whose residual exceeds $10^{-7}$ even when \code{success} is
\code{True} (\file{\_mean\_field\_dae.py:492-494}) --- scipy sometimes reports
success with a residual above its own internal tolerance.
\end{gotcha}

\subsection{Dedup and the deterministic sort}

Many seeds converge to the same root, so the raw list is clustered:
\code{\_dedup\_roots} (\file{\_mean\_field\_dae.py:368}) merges two roots when
$\max|r - \text{kept}| < \code{atol} + \code{rtol}\cdot\text{scale}$, keeping one
representative per cluster. The deduplicated roots are then sorted ascending by
the \emph{first declared physical field's first population index}:

\begin{lstlisting}
# pipeline/_mean_field_dae.py:506-508
first_var = state_vars[0][0]
first_start = slices[first_var][0]
deduped.sort(key=lambda r: float(r[first_start]))
\end{lstlisting}

This gives a deterministic, reproducible labeling: \code{fixed\_point\_index=0}
always means ``the root with the smallest first-field value.'' The list of state
variables comes from \code{\_state\_variables} (\file{\_mean\_field\_dae.py:75}),
which reads them in declaration order and prefers each field's user-facing
\code{natural\_name} (\code{'n'}, \code{'v'}) over its internal MSR name
(\code{'dn'}, \code{'dv'}).

\begin{gotcha}
The sort key is fragile by construction. Reordering the \code{physical\_field}
declarations in the theory changes which field is ``first,'' and therefore which
root \code{fixed\_point\_index=0} selects. If a previously-correct theory starts
picking the wrong branch after an edit, check whether the field declaration order
moved.
\end{gotcha}

\section{Linear stability}

Not every fixed point is a legitimate place to expand. The loop expansion is only
well-defined about a \emph{linearly stable} saddle: perturbations must decay,
otherwise the free propagator carries a growing mode and the perturbative series
diverges. Stability is read from the linearization of the DAE about the root,
implemented in \code{linear\_stability} (\file{\_mean\_field\_dae.py:631}).

\subsection{The matrix pencil}

Linearize $F_k(x,\code{Dt}) = 0$ about $x = x^{*}$, writing a perturbation
$\delta x \propto \ee^{\sigma t}$ so that \code{Dt} acting on it becomes
multiplication by $\sigma$. To first order,
\begin{equation}
  \sum_j \Big[\,\frac{\partial F_k}{\partial x_j}\Big|_{x^{*},\,\code{Dt}=0}
    + \sigma\,\frac{\partial^{2} F_k}{\partial \code{Dt}\,\partial x_j}\Big|_{x^{*}}
    \Big]\,\delta x_j = 0.
\end{equation}
Defining two real $M\times M$ matrices ($M$ = total number of scalar state
variables),
\begin{align}
  B_{kj} &= \frac{\partial F_k}{\partial x_j}\bigg|_{x^{*},\,\code{Dt}=0}
    && \text{(the algebraic Jacobian)}, \\
  A_{kj} &= \frac{\partial^{2} F_k}{\partial \code{Dt}\,\partial x_j}\bigg|_{x^{*}}
    && \text{(the ``mass'' matrix, the coefficient of $\sigma$)},
\end{align}
the linearized condition is $(\sigma A + B)\,\delta x = 0$ --- a
\term{generalized eigenvalue problem}. Rearranged,
$(-B)\,\delta x = \sigma A\,\delta x$, so the stability exponents $\sigma$ are the
eigenvalues of the matrix pencil $(-B, A)$.

\begin{defn}
A \textbf{generalized eigenvalue problem} is $C\,v = \lambda D\,v$ for two
matrices $C, D$. It reduces to the ordinary eigenproblem when $D$ is the identity.
It is the \emph{correct} tool here because the mass matrix $A$ is \emph{singular}:
algebraic equations (those without \code{Dt}) have zero rows in $A$, so $A$ cannot
be inverted and one cannot simply take the eigenvalues of $A^{-1}B$.
\end{defn}

The matrices are built by symbolic differentiation at the root. Note that the
\code{Dt} derivative for $A$ is taken \emph{first}, which drops the polynomial
degree before differentiating with respect to $x_j$:

\begin{lstlisting}
# pipeline/_mean_field_dae.py:798-811
for k, F_k in enumerate(residuals_sym):
    # B[k, j] = dF_k/dx_j at Dt=0, x=x*.
    F_k_alg = F_k.subs({Dt_sym: 0})
    for j, x_j in enumerate(flat_syms):
        d = diff(F_k_alg, x_j)
        B_mat[k, j] = float(d.subs(root_subs))

    # A[k, j] = d^2 F_k / dDt dx_j at x*.  dDt first to drop the degree.
    F_k_dDt = diff(F_k, Dt_sym)
    for j, x_j in enumerate(flat_syms):
        d = diff(F_k_dDt, x_j)
        A_mat[k, j] = float(d.subs(root_subs).subs({Dt_sym: 0}))
\end{lstlisting}

The symbolic residuals \code{residuals\_sym} are themselves built by \code{eval}ing
the equation text in a \emph{Sage} namespace (Section~\ref{sec:mf-tools} contrasts
this with the NumPy namespace the residual path uses), wrapping single-position
state-variable symbol lists in \code{\_FieldScalar}
(\file{theory\_compiler.py:268}) so that bare-name arithmetic such as \code{xt*x}
works alongside indexed access \code{xt[0]*x[0]}.

\subsection{Solving the pencil and the stability verdict}

The generalized eigenproblem is solved with \code{scipy.linalg.eig} called with
two matrix arguments. The infinite/NaN eigenvalues --- the algebraic-constraint
modes from the zero rows of $A$ --- are filtered out, and the verdict is read from
the surviving finite eigenvalues:

\begin{lstlisting}
# pipeline/_mean_field_dae.py:817-832  (abridged)
sigmas, eigvecs = scipy.linalg.eig(-B_mat, A_mat)
finite_mask = np.isfinite(sigmas)
sigmas_finite = sigmas[finite_mask]
EIG_TOL = 1e-9
stable = bool(sigmas_finite.size > 0 and np.all(np.real(sigmas_finite) < -EIG_TOL))
unstable = [complex(s) for s in sigmas_finite if np.real(s) >= -EIG_TOL]
\end{lstlisting}

The convention is: the fixed point is stable iff there is at least one finite
eigenvalue \emph{and every} finite eigenvalue has $\mathrm{Re}(\sigma) < 0$. The
small tolerance \code{EIG\_TOL} $= 10^{-9}$ keeps an eigenvalue that is zero to
roundoff from being misclassified. A growing mode ($\mathrm{Re}\,\sigma \ge 0$)
marks the root unstable.

\begin{note}
The lazy import of Sage at \file{\_mean\_field\_dae.py:676} is deliberate. The DAE
\emph{root-finding} path is pure NumPy/SciPy and never imports Sage; \emph{only}
the stability classification needs symbolic differentiation, so Sage (a heavy
import) is loaded only when \code{linear\_stability} actually runs.
\end{note}

\subsection{The \code{stability\_analysis} toggle}

Whether the eigenvalue machinery runs at all is gated by
\code{model['stability\_analysis']} (default \code{False}), set on the builder via
\code{.stability\_analysis(...)} (\file{pipeline/theory.py:1187}):

\begin{description}
  \item[OFF (the default).] Skip stability entirely. Every converged root is
        selectable, and \code{fixed\_point\_index} ranges over all of them. Each
        record is tagged \code{stable=None} (to distinguish ``not classified''
        from ``classified unstable''). This is the \emph{correct} choice when
        every equation is algebraic --- voltages integrated out, no \code{Dt}
        anywhere --- because then $A \equiv 0$, every eigenvalue is infinite, and
        ``linear stability'' is meaningless. A Hawkes-type rate model is the
        canonical example.
  \item[ON.] Classify every root, filter to the linearly-stable subset, expose
        \code{fixed\_point\_index} over \emph{that} subset, and \emph{raise} if no
        stable root exists. Required for bistable or multi-saddle theories where
        you must pick the stable branch.
\end{description}

\begin{lstlisting}
# pipeline/_mean_field_dae.py:569-586  (abridged)
if stab_on:
    selectable_records = [r for r in all_root_records if r['stable']]
    if not selectable_records:
        raise ValueError(
            "... found N fixed point(s) but NONE were classified linearly "
            "stable.  ... (If your equations are all algebraic -- voltages "
            "integrated out, no Dt anywhere -- set ``.stability_analysis(False)`` "
            "on the TheoryBuilder; ``fixed_point_index`` will then range over "
            "every converged root.)")
\end{lstlisting}

\begin{gotcha}
Turning \code{stability\_analysis} ON for an all-algebraic theory is a guaranteed
failure: with no \code{Dt} anywhere, $A \equiv 0$, every eigenvalue is infinite,
the stable subset is empty, and the solver raises ``NONE were classified linearly
stable.'' The error message itself tells you the fix --- call
\code{.stability\_analysis(False)}.
\end{gotcha}

\begin{gotcha}
Per-root \code{linear\_stability} calls are wrapped in \code{except Exception}
(\file{\_mean\_field\_dae.py:542-552}): a stability failure marks that root
\emph{unstable} (with an \code{'error'} annotation) rather than aborting the whole
solve. The consequence is that a bug in the stability code can make a genuinely
stable root silently vanish from the selectable set.
\end{gotcha}

\subsection{Index selection and clamping}

Finally, \code{fixed\_point\_index} selects from the selectable subset. An
out-of-range index is \emph{clamped} with a warning rather than raising:

\begin{lstlisting}
# pipeline/_mean_field_dae.py:594-603
if fixed_point_index < 0 or fixed_point_index >= n_selectable:
    clamped = max(0, min(fixed_point_index, n_selectable - 1))
    warnings.warn(
        f"... requested fixed_point_index={fixed_point_index} but only "
        f"{n_selectable} {index_label} found ...; using index {clamped}.",
        stacklevel=2)
    fixed_point_index = clamped
\end{lstlisting}

\begin{gotcha}
An out-of-range \code{fixed\_point\_index} never errors --- it clamps to the valid
range, and a single \code{warnings.warn} is the only signal. In a noisy notebook
that warning is easy to miss, so if you asked for the third fixed point and got
the first, check for a clamp warning.
\end{gotcha}

\section{From saddle to diagrams}

The handoff to the perturbative expansion is \emph{purely substitutional}. The
propagator kernel $K_{\mathrm{ft}}(\omega)$ and the vertices are symbolic Sage
expressions containing parameter symbols and saddle symbols. The mean-field
solver's deliverable, \code{num\_params}, is a dict
$\{\text{SR symbol}\to\text{float}\}$ covering every numeric parameter (\code{E1},
\code{w12}, \code{tau}, \dots), every saddle value (\code{nstar[i]},
\code{vstar[i]}, \code{mstar[i]}), every transfer-function derivative at the
saddle (\code{phi1\_i}, \code{phi2\_i}, \dots, on the legacy path), and the
\code{phi0\_i} $\to$ \code{nstar\_i} identities.

Downstream, \code{compute\_poles\_and\_residues(prop, num\_params)}
(\file{pipeline/\_propagator.py:851}) substitutes \code{num\_params} into
$K_{\mathrm{ft}}$, obtaining a numeric matrix-valued rational function of
$\omega$. It then finds the retarded poles $\omega_k$ (those with
$\mathrm{Im}\,\omega > 0$) and the residue matrices that assemble the free
propagator $G(\tau)$ (the sibling chapter, Chapter~\ref{ch:propagator}, covers
that). From there the loop corrections build the connected cumulant $C(\tau)$. The
saddle enters as \emph{the numbers that turn symbolic diagrams into numeric ones}
--- there is no further re-solve; the saddle is frozen and the loops are
corrections about it.

\begin{note}
The symbolic-derivative renaming that \emph{creates} the \code{phi1\_i},
\code{phi2\_i} symbols lives upstream in the compiler, in
\code{ns.\_deriv\_rename\_subs} (\file{pipeline/theory\_compiler.py:1048}): the
compiler differentiates each formal function at its saddle and registers a rename
$\langle\text{derivative}\rangle \to \code{SR.var('phik\_<i+1>')}$. The legacy
solver then supplies the \emph{numeric} value of each \code{phik\_i}. The DAE
compatibility wrapper does \emph{not} populate \code{phi\_deriv\_vals} (it returns
\code{\{\}}); the propagator builds those derivatives symbolically in that path.
\end{note}

\subsection{The compatibility wrapper}

\code{solve\_mean\_field\_dae\_compat} (\file{\_mean\_field\_dae.py:853}) is what
\code{compute.py} actually calls on the DAE path. It is a drop-in replacement for
the legacy \code{solve\_mean\_field}: it runs \code{solve\_mean\_field\_dae},
builds \code{param\_subs} in the legacy Sage-symbol convention, and bakes the
selected root's saddle values into \code{num\_params} by looking up the Sage
arrays on \code{ft.\_ns}:

\begin{lstlisting}
# pipeline/_mean_field_dae.py:943-950
num_params = dict(param_subs)
for saddle_name, vals in mf_values.items():
    sr_array = getattr(ft._ns, saddle_name, None)
    if sr_array is None:
        continue
    for i, v in enumerate(vals):
        num_params[sr_array[i]] = float(v)
\end{lstlisting}

It returns the legacy-shape keys (\code{nstar\_vals}, \code{vstar\_vals},
\code{phi\_deriv\_vals=\{\}}, \code{num\_params}, \code{param\_subs},
\code{saddle\_values}) \emph{plus} the DAE extras (\code{mf\_all\_roots},
\code{mf\_stable\_roots}, \code{mf\_unstable\_roots}, \code{mf\_index\_used},
\code{state\_var\_order}, \code{n\_seeds\_converged}). \code{compute\_cumulants}
surfaces those extras onto its returned result dict, so a notebook can introspect
every fixed point and its stability --- which makes them the natural first stop
when a DAE saddle behaves unexpectedly.

\section{External tooling primer}\label{sec:mf-tools}

This subsystem leans on four external pieces. Since the manual assumes you know
the physics but not necessarily the tooling, here is each from scratch, with the
exact import lines.

\subsection{SageMath and the symbolic ring \code{SR}}

\term{SageMath} is a large open-source mathematics system built on top of Python;
it bundles many computer-algebra packages behind one Python API. The piece this
subsystem uses is its \emph{symbolic ring} \code{SR} (the ``Symbolic Ring''),
which lets you build, differentiate, and substitute into expressions made of
symbols (\code{x}, \code{E1}, \code{nstar}) rather than numbers. Think of it as a
calculator that does algebra with letters.

\begin{itemize}
  \item \code{SR.var('name')} returns (and globally registers) a symbolic
        variable; calling it twice with the same name returns the \emph{same}
        symbol. This is why a substitution dict keyed by \code{SR.var('E1')}
        matches the \code{E1} the propagator built separately.
  \item \code{SR(value)} coerces a Python number or another expression into the
        ring.
  \item \code{diff(expr, var, k)} computes the $k$-th symbolic derivative.
  \item \code{.subs(dict)} substitutes symbols $\to$ values (or expressions);
        chained \code{.subs(...).subs(...)} applies left to right. After all
        symbols are numbers, \code{float(expr)} evaluates to a Python float.
\end{itemize}

The legacy solver imports Sage at module top (\file{\_mean\_field.py:10}:
\code{from sage.all import SR, diff}). The DAE solver imports it \emph{lazily},
inside \code{linear\_stability} (\file{\_mean\_field\_dae.py:676-681}), because
only stability needs symbolic differentiation.

\begin{note}
The contrast between the two math namespaces is the key to the DAE design. The
\emph{residual} path binds \code{tanh} $\to$ \code{np.tanh} (\code{\_MATH\_NS}),
so an equation evaluates to a \emph{float}. The \emph{stability} path binds
\code{tanh} $\to$ Sage's \code{tanh} (\code{sage\_math\_ns}), so the same equation
stays \emph{symbolic} and can be differentiated. The very same text is evaluated
two different ways for two different purposes.
\end{note}

\subsection{SciPy --- numerical solvers}

\term{SciPy} is the standard scientific-computing library, built on NumPy. It
provides numerical (not symbolic) algorithms; nothing in it knows about symbols.

\begin{itemize}
  \item \code{scipy.optimize.fsolve} (legacy solver, \file{\_mean\_field.py:11})
        finds a root of $f(x)=0$ from a starting guess using MINPACK's modified
        Powell hybrid method. With \code{full\_output=True} it returns a 4-tuple
        \code{(x, infodict, ier, mesg)}; \code{ier == 1} signals clean
        convergence.
  \item \code{scipy.optimize.root} (DAE solver, imported as \code{\_scipy\_root}
        at \file{\_mean\_field\_dae.py:27}) is the newer, more general
        root-finder. With \code{method='hybr'} it is the same Powell hybrid, but
        it returns a richer \code{OptimizeResult} with \code{.x} and
        \code{.success}.
  \item \code{scipy.linalg.eig} (stability, \file{\_mean\_field\_dae.py:682})
        solves eigenvalue problems. Called with \emph{two} matrices,
        \code{eig(C, D)}, it solves the generalized problem $C v = \lambda D v$.
        The code calls \code{eig(-B, A)}, so the returned eigenvalues \emph{are}
        the stability exponents $\sigma$.
\end{itemize}

\subsection{NumPy --- numerical arrays}

\term{NumPy} is the foundational array library: the \code{ndarray}, elementwise
math (\code{np.tanh}, \code{np.exp}), and reductions (\code{np.max}, \code{np.abs},
\code{np.all}, \code{np.isfinite}). It is the substrate SciPy stands on. The DAE
solver uses it to coerce parameters into arrays so that \code{w[i,j]} (2-D) and
\code{Em[i]} (1-D) work under \code{eval} (\code{\_numpy\_params},
\file{\_mean\_field\_dae.py:198}), to draw the random seeds via the modern
generator \code{np.random.default\_rng} (\file{\_mean\_field\_dae.py:475}), and to
pack and filter the residual and root vectors.

\subsection{\code{eval} and the namespace conventions}

\code{eval(source, globals\_dict)} executes a Python expression string in a
supplied namespace and returns its value. The DAE solver \code{eval}s the
(Sage-to-Python translated) equation text directly. To keep that safe and
predictable it passes a \emph{closed} namespace with
\code{'\_\_builtins\_\_': \{\}} (so no builtins leak in) and binds exactly the
names an expression may reference. The two non-obvious rules --- pass the namespace
as \emph{globals} (the PEP-3104 comprehension-scoping fix) and the empty-builtins
sandbox --- were covered in the gotchas above; they are the difference between an
equation that evaluates and one that raises a swallowed \code{NameError}.

\section{Data structures}

For reference, the dictionaries that cross this subsystem's boundaries.

\subsection{The legacy solver return dict}

\begin{lstlisting}
{
  'nstar_vals':     [float, ...],       # rate saddles (per pop / concatenated)
  'vstar_vals':     [float, ...],       # voltage saddles
  'phi_deriv_vals': {(dk, i): float},   # phi^(dk)(v*_i); {} on the hetero path
  'num_params':     {SR var: float},    # THE handoff dict
  'param_subs':     {SR var: float},    # raw user params only
  # hetero path also adds:
  'saddle_values':  {saddle_name: [float]},
  'saddle_info':    {saddle_name: {'pop', 'size', 'sr_array'}},
}
\end{lstlisting}

\subsection{The DAE solver return dict}

\begin{lstlisting}
{
  'mf_values':         {'<var>star': [floats]},   # the selected root
  'mf_all_roots':      [ {'values': {...},         # every root, sorted, annotated
                          'stable': bool|None,
                          'eigenvalues_finite': np.ndarray[complex]}, ... ],
  'mf_stable_roots':   [ {'<var>star': [...]}, ... ],  # selectable subset
  'mf_unstable_roots': [ {...}, ... ],             # only meaningful when stab ON
  'mf_index_used':     int,                        # index after clamping
  'state_var_order':   [var_name, ...],            # declaration order (sort key)
  'n_seeds_converged': int,                        # pre-dedup convergence count
  'stability_analysis': bool,                      # was classification run?
}
\end{lstlisting}

\subsection{The \code{linear\_stability} return dict}

\begin{lstlisting}
{
  'stable':               bool,
  'eigenvalues_finite':   np.ndarray[complex],   # dynamical modes sigma
  'eigenvalues_all':      np.ndarray[complex],   # includes +/-inf (constraints)
  'A':                    np.ndarray[float] (MxM),  # d^2 F / dDt dx   (pencil)
  'B':                    np.ndarray[float] (MxM),  # dF / dx at Dt=0  (Jacobian)
  'unstable_eigenvalues': [complex, ...],        # finite sigma with Re >= 0
}
\end{lstlisting}

\section{Pitfalls and diagnostics}

A consolidated checklist of the failure modes scattered through this chapter, for
when a saddle misbehaves:

\begin{itemize}
  \item \textbf{Two solvers, one switch.} Which solver runs is purely
        \code{model.get('equations')} truthiness. A theory carrying both
        \code{.equation(...)} and \code{mf\_bg\_conditions} silently takes the DAE
        path.
  \item \textbf{Singularities are penalized, not raised.} Both solvers turn
        symbolic singularities (spike-reset $1/(1-n^{*})$) into a $10^{10}$
        residual to let the solver back away. But a root that \emph{converges}
        into a near-pole is caught by explicit finiteness checks, which raise a
        \code{RuntimeError} with guidance (\file{\_mean\_field.py:193-217}).
  \item \textbf{\code{phi0\_i} must be resolved or the correlator is zero.}
        Template-built theories keep $\phi$ formal; the noise coefficient
        $-\tfrac12\,\code{phi0\_i}$ stays symbolic unless \code{phi0\_i} $\to$
        \code{nstar\_i} is baked in (\file{\_mean\_field.py:260-279}).
  \item \textbf{DAE \code{eval} namespace must be globals, with empty builtins.}
        Get the globals/locals wrong and comprehensions raise \code{NameError};
        reference an unbound builtin and the exception is swallowed into a
        $10^{10}$ penalty that looks like non-convergence.
  \item \textbf{\code{stability\_analysis} must be OFF for all-algebraic
        theories.} Otherwise $A\equiv 0$, no eigenvalue is finite, and the solver
        raises ``NONE were classified linearly stable.''
  \item \textbf{\code{fixed\_point\_index} clamps with a warning, never errors.}
        An out-of-range index is silently clamped.
  \item \textbf{Sort-key fragility.} Roots are ordered by the first declared
        field's first index; reordering field declarations changes which root is
        index $0$.
  \item \textbf{Swallowed per-root stability failures.} A bug in
        \code{linear\_stability} can make a genuinely stable root vanish from the
        selectable set, because the per-root call is wrapped in
        \code{except Exception}.
  \item \textbf{Reading \code{mf\_all\_roots}.} When a DAE theory picks an
        unexpected saddle, inspect \code{res['mf\_all\_roots']} (every root with
        its stability flag and finite eigenvalues) and call
        \code{linear\_stability(model, fundamental, root)} directly on a root to
        see its $A$, $B$, and eigenvalues.
\end{itemize}

\begin{note}
\code{Laplacian} and \code{Dt} are both hard-substituted to $0$ on the
homogeneous mean field. Stability to \emph{spatial} ($k\neq 0$) perturbations ---
the pattern-formation question --- is \emph{not} handled here; it lives in the
propagator's momentum dependence, as the source notes at
\file{\_mean\_field\_dae.py:769-773}.
\end{note}
